% Verbatim from minor_report/chapters/chapter_4_results.tex lines 1730-1993 (retired report),
% headings demoted by 0 level(s). Do not edit here; edit the source of the numbers.

\subsection{Where each pipeline is weak, block by block}
\label{sec:block-performance}

\Cref{fig:blockperf-static,fig:blockperf-dynamic} put every instrumented stage of
the three pipelines on one page and colour it by a stated criterion.
\Cref{tab:blockperf-static-orb} through \cref{tab:blockperf-dynamic-vins} give the evidence behind every
colour: the metric, its value, the criterion it was tested against, and which of
three rules produced the verdict. Both the figures' fills and the tables are
generated by \repo{script/block_performance.py} from the run artifacts; no colour
in either figure was chosen by hand.

The three rules are deliberately distinguished, because they answer different
questions and only one of them is about accuracy:

\begin{description}[leftmargin=2.2em,style=nextline]
  \item[ablation] The block was removed, disabled or swapped and the run repeated.
    The value is the ratio of aligned ATE without it to aligned ATE with it. This
    is the only rule that speaks to trajectory error.
  \item[latency] The block's median time as a share of its stage, or a stage
    against the \SI{66.0}{\milli\second} inter-frame budget implied by
    \SI{30.3}{\hertz} imagery at $0.5\times$ replay. Describes cost, not accuracy.
  \item[outcome] A counter the block emits, against a stated threshold---accepted
    loop candidates, usable solver solutions, filter consistency.
\end{description}

\paragraph{Where the criteria come from.}
Every verdict is a comparison against a threshold, and those thresholds do not all
have the same standing. The tables label each one, and the distinction matters
more than the values:

\begin{description}[leftmargin=2.4em,style=nextline]
  \item[\textit{derived}] Follows from something outside this report. There are
    two. The \SI{66.0}{\milli\second} frame budget is arithmetic on the recorded
    conditions, $1000/(30.3\,\mathrm{Hz}\times0.5$ replay$)$: a per-frame stage
    slower than this cannot keep up, and no judgement is involved. The consistency
    target $\mathrm{ANEES}=3.0$ is the standard result that a correctly scaled
    covariance makes the normalised error squared average the degrees of freedom,
    which for a pose error is three.
  \item[\textit{convention}] A round number adopted here so the rule is stated
    rather than left implicit. Nothing external fixes \SI{40}{\percent} of a stage
    rather than \SI{35}{\percent}, $2.0\times$ for a load-bearing block rather than
    $1.5\times$, \SI{10}{\percent} of loop candidates accepted, or
    \SI{99}{\percent} of solver calls returning a usable solution. These are
    contestable, which is why the measured value sits beside the threshold in
    every row: a reader who prefers a different line can redraw it without
    re-running anything.
\end{description}

\begin{queried}[Withdrawn criterion]
An earlier revision of these tables graded VINS-Fusion's solver quality against
``median final cost $>5000$''. That threshold was chosen after observing 1176 in
static scenes and 14110 in dynamic ones, so it separated two numbers already in
hand and the \flag{``poor'' verdict it produced was guaranteed by the choice
rather than tested by it}. Raw Ceres cost also has no absolute scale, depending on
residual count and weighting, so no fixed threshold on it is meaningful in
principle. It is replaced by the solver's own \repo{solver_solution_usable} flag,
which is binary and emitted by the solver rather than invented here:
\SI{99.64}{\percent} of solves in static scenes and \SI{99.55}{\percent} in
dynamic ones return a usable solution, and the block is now graded \perfmark{good}
in both. The non-convergence rate---\SI{22.9}{\percent} of solves in static scenes
and \SI{45.3}{\percent} in dynamic ones reaching the iteration or
\SI{80}{\milli\second} cap---is reported in the note column and deliberately not
graded, because any threshold on it would repeat the same error.
\end{queried}

\begin{queried}[Reading the colours]
A red block is \emph{not} a claim that it caused the trajectory error, unless its
rule is \texttt{ablation}. Only four blocks in the entire project have an
ablation behind them: the EKF's wheel correction, its IMU prediction, its
yaw-rate model, and ORB-SLAM3's loop closing. Every other red marks a block that
is expensive or that emits a poor counter, which is a different and weaker
statement. Yellow means no claim at all---not instrumented, not exercised in these
runs, or not implemented in this fork.
\end{queried}

\begin{figure}[p]
  \centering
  %% Generated by script/block_performance.py -- do not edit.
\expandafter\def\csname blockcolour@static@orbextract\endcsname{perfpoor}
\expandafter\def\csname blockcolour@static@orbimupre\endcsname{perfgood}
\expandafter\def\csname blockcolour@static@orbposepred\endcsname{perfgood}
\expandafter\def\csname blockcolour@static@orbtracklocal\endcsname{perfgood}
\expandafter\def\csname blockcolour@static@orbkfdec\endcsname{perfgood}
\expandafter\def\csname blockcolour@static@orbkfins\endcsname{perfgood}
\expandafter\def\csname blockcolour@static@orbmpcull\endcsname{perfgood}
\expandafter\def\csname blockcolour@static@orbnewpoints\endcsname{perfgood}
\expandafter\def\csname blockcolour@static@orblba\endcsname{perfpoor}
\expandafter\def\csname blockcolour@static@orbkfcull\endcsname{perfgood}
\expandafter\def\csname blockcolour@static@orbtrackingtotal\endcsname{perfgood}
\expandafter\def\csname blockcolour@static@orbimuinit\endcsname{perfgood}
\expandafter\def\csname blockcolour@static@orbimuref\endcsname{perfgood}
\expandafter\def\csname blockcolour@static@orbplacerec\endcsname{perfpoor}
\expandafter\def\csname blockcolour@static@orbloopfusion\endcsname{perfpoor}
\expandafter\def\csname blockcolour@static@orbfullba\endcsname{perfnone}
\expandafter\def\csname blockcolour@static@orbyolo\endcsname{perfnone}
\expandafter\def\csname blockcolour@static@vinstrack\endcsname{perfgood}
\expandafter\def\csname blockcolour@static@vinsqueue\endcsname{perfpoor}
\expandafter\def\csname blockcolour@static@vinsimuprop\endcsname{perfgood}
\expandafter\def\csname blockcolour@static@vinstri\endcsname{perfgood}
\expandafter\def\csname blockcolour@static@vinsconstruct\endcsname{perfgood}
\expandafter\def\csname blockcolour@static@vinssolve\endcsname{perfpoor}
\expandafter\def\csname blockcolour@static@vinsmarg\endcsname{perfgood}
\expandafter\def\csname blockcolour@static@vinsoutlier\endcsname{perfgood}
\expandafter\def\csname blockcolour@static@vinsslide\endcsname{perfgood}
\expandafter\def\csname blockcolour@static@vinsbackendtotal\endcsname{perfpoor}
\expandafter\def\csname blockcolour@static@vinssolverq\endcsname{perfgood}
\expandafter\def\csname blockcolour@static@vinsloop\endcsname{perfnone}
\expandafter\def\csname blockcolour@static@vinsrtab\endcsname{perfnone}
\expandafter\def\csname blockcolour@static@vinsyolo\endcsname{perfnone}
\expandafter\def\csname blockcolour@static@ekfwheelupd\endcsname{perfgood}
\expandafter\def\csname blockcolour@static@ekfbiasobs\endcsname{perfgood}
\expandafter\def\csname blockcolour@static@ekfimupred\endcsname{perfgood}
\expandafter\def\csname blockcolour@static@ekfyaw\endcsname{perfpoor}
\expandafter\def\csname blockcolour@static@ekfanchor\endcsname{perfnone}
\expandafter\def\csname blockcolour@static@ekfcov\endcsname{perfpoor}
%
  \resizebox{\textwidth}{!}{\pipelinefig{static}}
  \caption[Per-block performance in the \textbf{static} condition, over the five \repo{allsensor} datasets.]{Per-block performance in the \textbf{static} condition, over the five
  \repo{allsensor} datasets. The EKF baseline appears here and not in
  \cref{fig:blockperf-dynamic} because no dynamic dataset carries wheel encoders.
  Evidence for every colour is in \cref{tab:blockperf-static-orb,tab:blockperf-static-vins,tab:blockperf-static-ekf}, one table per algorithm.}
  \label{fig:blockperf-static}
\end{figure}
\clearpage

\begin{table}[p]
\centering
\caption[Per-block evidence for \textbf{ORB-SLAM3} in the static condition, \cref{fig:blockperf-static}.]{Per-block evidence for \textbf{ORB-SLAM3} in the static condition, \cref{fig:blockperf-static}. Five \repo{allsensor} datasets, median across runs, loop closing enabled (\repo{logging_20260916_01}).}
\label{tab:blockperf-static-orb}
\scriptsize
\setlength{\tabcolsep}{3pt}
%% Generated by script/block_performance.py -- do not edit.
\begin{tabularx}{\textwidth}{>{\raggedright\arraybackslash}p{0.155\textwidth}l>{\raggedright\arraybackslash}X r>{\raggedright\arraybackslash}p{0.20\textwidth}l}
\toprule
Block & Rule & Metric & Value & Criterion / note & Verdict \\
\midrule
\multicolumn{6}{l}{\textbf{Tracking}} \\
\addlinespace[2pt]
\quad Extract ORB (+preproc, stereo) & latency & median image\_preprocessing\_feature\_stereo\_ms & \num{33.1} & \textit{convention:} $<$40\% of stage (61.4\% of stage) & \perfmark{poor} \\
\addlinespace[3.5pt]
\quad IMU preintegration & latency & median imu\_preintegration\_ms & \num{0.120} & \textit{convention:} $<$40\% of stage (0.2\% of stage) & \perfmark{good} \\
\addlinespace[3.5pt]
\quad Initial pose estimation & latency & median pose\_prediction\_ms & \num{0.087} & \textit{convention:} $<$40\% of stage (0.2\% of stage) & \perfmark{good} \\
\addlinespace[3.5pt]
\quad Track local map & latency & median local\_map\_tracking\_ms & \num{16.5} & \textit{convention:} $<$40\% of stage (30.7\% of stage) & \perfmark{good} \\
\addlinespace[3.5pt]
\quad New keyframe decision & latency & median keyframe\_decision\_ms & \num{0.169} & \textit{convention:} $<$40\% of stage (0.3\% of stage) & \perfmark{good} \\
\addlinespace[3.5pt]
\quad TRACKING total & latency & median tracking\_total\_ms & \num{53.8} & \textit{derived:} $<$\SI{66.0}{\milli\second} frame budget & \perfmark{good} \\
\addlinespace[3.5pt]
\quad Dynamic feature removal (YOLO) & none & -- & -- & -- (filter disabled in every run reported here) & \perfmark{none} \\
\addlinespace[3.5pt]
\addlinespace[5pt]
\multicolumn{6}{l}{\textbf{Local mapping}} \\
\addlinespace[2pt]
\quad KeyFrame insertion & latency & median keyframe\_insertion\_ms & \num{10.3} & \textit{convention:} $<$40\% of stage (3.3\% of stage) & \perfmark{good} \\
\addlinespace[3.5pt]
\quad Recent MapPoints culling & latency & median map\_point\_culling\_ms & \num{0.383} & \textit{convention:} $<$40\% of stage (0.1\% of stage) & \perfmark{good} \\
\addlinespace[3.5pt]
\quad New points creation & latency & median map\_point\_creation\_fusion\_ms & \num{41.5} & \textit{convention:} $<$40\% of stage (13.3\% of stage) & \perfmark{good} \\
\addlinespace[3.5pt]
\quad Local BA & latency & median local\_ba\_ms & \num{244.4} & \textit{convention:} $<$40\% of stage (78.1\% of stage) & \perfmark{poor} \\
\addlinespace[3.5pt]
\quad Local keyframes culling & latency & median keyframe\_culling\_ms & \num{10.9} & \textit{convention:} $<$40\% of stage (3.5\% of stage) & \perfmark{good} \\
\addlinespace[3.5pt]
\quad IMU initialization & latency & median imu\_initialization\_ms & \num{0.002} & \textit{convention:} $<$40\% of stage (negligible) & \perfmark{good} \\
\addlinespace[3.5pt]
\quad IMU scale refinement & latency & median inertial\_refinement\_ms & \num{0.001} & \textit{convention:} $<$40\% of stage (negligible) & \perfmark{good} \\
\addlinespace[3.5pt]
\addlinespace[5pt]
\multicolumn{6}{l}{\textbf{Loop closing}} \\
\addlinespace[2pt]
\quad Place recognition & outcome & candidates accepted & 4 of 281 & \textit{convention:} $>$\SI{10}{\percent} of candidates accepted (2494 events) & \perfmark{poor} \\
\addlinespace[3.5pt]
\quad Loop fusion + essential graph & ablation & ATE(on)/ATE(off) & 1.37 & \textit{convention:} $<$1.0 (enabling it should help) (one run per configuration) & \perfmark{poor} \\
\addlinespace[3.5pt]
\quad Full BA & outcome & global BA executions & -- & -- (see \cref{tab:orb-loop}) & \perfmark{none} \\
\addlinespace[3.5pt]
\bottomrule
\end{tabularx}

\end{table}
\clearpage
\begin{table}[p]
\centering
\caption[Per-block evidence for \textbf{VINS-Fusion} in the static condition, \cref{fig:blockperf-static}.]{Per-block evidence for \textbf{VINS-Fusion} in the static condition, \cref{fig:blockperf-static}. Five \repo{allsensor} datasets, median across runs.}
\label{tab:blockperf-static-vins}
\scriptsize
\setlength{\tabcolsep}{3pt}
%% Generated by script/block_performance.py -- do not edit.
\begin{tabularx}{\textwidth}{>{\raggedright\arraybackslash}p{0.155\textwidth}l>{\raggedright\arraybackslash}X r>{\raggedright\arraybackslash}p{0.20\textwidth}l}
\toprule
Block & Rule & Metric & Value & Criterion / note & Verdict \\
\midrule
\multicolumn{6}{l}{\textbf{Front end}} \\
\addlinespace[2pt]
\quad Feature tracking & latency & median feature\_tracking\_ms & \num{55.2} & $<$\SI{66.0}{\milli\second} frame budget & \perfmark{good} \\
\addlinespace[3.5pt]
\quad Feature queue & latency & median feature\_queue\_wait\_ms & \num{403.7} & $<$\SI{66.0}{\milli\second} frame budget (backlog 4.0 frames) & \perfmark{poor} \\
\addlinespace[3.5pt]
\quad Persistent-feature removal (YOLO) & none & -- & -- & -- (filter disabled in every run reported here) & \perfmark{none} \\
\addlinespace[3.5pt]
\addlinespace[5pt]
\multicolumn{6}{l}{\textbf{Back end}} \\
\addlinespace[2pt]
\quad IMU propagation & latency & median imu\_propagation\_ms & \num{0.272} & \textit{convention:} $<$40\% of stage (0.4\% of stage) & \perfmark{good} \\
\addlinespace[3.5pt]
\quad Triangulation & latency & median triangulation\_ms & \num{0.063} & \textit{convention:} $<$40\% of stage (0.1\% of stage) & \perfmark{good} \\
\addlinespace[3.5pt]
\quad Problem construction & latency & median problem\_construction\_ms & \num{3.2} & \textit{convention:} $<$40\% of stage (4.7\% of stage) & \perfmark{good} \\
\addlinespace[3.5pt]
\quad Ceres solve & latency & median solver\_ms & \num{48.3} & \textit{convention:} $<$40\% of stage (71.9\% of stage) & \perfmark{poor} \\
\addlinespace[3.5pt]
\quad Marginalization & latency & median marginalization\_ms & \num{4.1} & \textit{convention:} $<$40\% of stage (6.1\% of stage) & \perfmark{good} \\
\addlinespace[3.5pt]
\quad Outlier rejection & latency & median outlier\_rejection\_ms & \num{0.155} & \textit{convention:} $<$40\% of stage (0.2\% of stage) & \perfmark{good} \\
\addlinespace[3.5pt]
\quad Slide window & latency & median slide\_window\_ms & \num{0.212} & \textit{convention:} $<$40\% of stage (0.3\% of stage) & \perfmark{good} \\
\addlinespace[3.5pt]
\quad BACK END total & latency & median backend\_total\_ms & \num{67.2} & $<$\SI{66.0}{\milli\second} frame budget & \perfmark{poor} \\
\addlinespace[3.5pt]
\quad Solver convergence & outcome & solves returning a usable solution & 99.64\% & \textit{convention:} $\geq$\SI{99}{\percent} usable (\SI{22.8}{\percent} hit the iteration or \SI{80}{\milli\second} cap without converging; reported, not graded) & \perfmark{good} \\
\addlinespace[3.5pt]
\addlinespace[5pt]
\multicolumn{6}{l}{\textbf{Loop closing}} \\
\addlinespace[2pt]
\quad Loop fusion & none & -- & -- & -- (not implemented in this fork) & \perfmark{none} \\
\addlinespace[3.5pt]
\quad RTAB-Map correction & none & -- & -- & -- (separate process; correction never injected) & \perfmark{none} \\
\addlinespace[3.5pt]
\bottomrule
\end{tabularx}

\end{table}
\clearpage
\begin{table}[p]
\centering
\caption[Per-block evidence for the \textbf{EKF wheel\,+\,IMU} baseline, \cref{fig:blockperf-static}.]{Per-block evidence for the \textbf{EKF wheel\,+\,IMU} baseline, \cref{fig:blockperf-static}. Every row but the last two is an ablation: the block was removed or swapped and all five datasets re-run, so these are the only verdicts in this section that speak to trajectory error.}
\label{tab:blockperf-static-ekf}
\scriptsize
\setlength{\tabcolsep}{3pt}
%% Generated by script/block_performance.py -- do not edit.
\begin{tabularx}{\textwidth}{>{\raggedright\arraybackslash}p{0.155\textwidth}l>{\raggedright\arraybackslash}X r>{\raggedright\arraybackslash}p{0.20\textwidth}l}
\toprule
Block & Rule & Metric & Value & Criterion / note & Verdict \\
\midrule
\multicolumn{6}{l}{\textbf{EKF}} \\
\addlinespace[2pt]
\quad Wheel speed correction & ablation & ATE without the wheel update / with & 54.08 & \textit{convention:} $\geq$2.0$\times$ (block is load-bearing) (median over 5 bags; joint with the other row -- one flag removes both) & \perfmark{good} \\
\addlinespace[3.5pt]
\quad Gyro-bias observation & ablation & ATE without the wheel update / with & 54.08 & \textit{convention:} $\geq$2.0$\times$ (block is load-bearing) (median over 5 bags; joint with the other row -- one flag removes both) & \perfmark{good} \\
\addlinespace[3.5pt]
\quad IMU prediction & ablation & ATE of wheel-only / EKF & 6.56 & \textit{convention:} $\geq$2.0$\times$ (block is load-bearing) (median over 5 bags) & \perfmark{good} \\
\addlinespace[3.5pt]
\quad Yaw-rate model (bicycle) & ablation & ATE with differential / with bicycle & 1.02 & \textit{convention:} $\geq$2.0$\times$ (block is load-bearing) (median over 5 bags) & \perfmark{poor} \\
\addlinespace[3.5pt]
\quad Initial pose anchor & none & -- & -- & -- (ground-truth pose; not a performance block) & \perfmark{none} \\
\addlinespace[3.5pt]
\quad Covariance propagation & outcome & ANEES (3 DOF) & -- & \textit{derived:} 3.0 for a consistent filter (64.7--123.2; see \cref{tab:proprio-covariance}) & \perfmark{poor} \\
\addlinespace[3.5pt]
\bottomrule
\end{tabularx}

\end{table}
\clearpage
\begin{figure}[p]
  \centering
  %% Generated by script/block_performance.py -- do not edit.
\expandafter\def\csname blockcolour@dynamic@orbextract\endcsname{perfpoor}
\expandafter\def\csname blockcolour@dynamic@orbimupre\endcsname{perfgood}
\expandafter\def\csname blockcolour@dynamic@orbposepred\endcsname{perfgood}
\expandafter\def\csname blockcolour@dynamic@orbtracklocal\endcsname{perfgood}
\expandafter\def\csname blockcolour@dynamic@orbkfdec\endcsname{perfgood}
\expandafter\def\csname blockcolour@dynamic@orbkfins\endcsname{perfgood}
\expandafter\def\csname blockcolour@dynamic@orbmpcull\endcsname{perfgood}
\expandafter\def\csname blockcolour@dynamic@orbnewpoints\endcsname{perfgood}
\expandafter\def\csname blockcolour@dynamic@orblba\endcsname{perfpoor}
\expandafter\def\csname blockcolour@dynamic@orbkfcull\endcsname{perfgood}
\expandafter\def\csname blockcolour@dynamic@orbtrackingtotal\endcsname{perfgood}
\expandafter\def\csname blockcolour@dynamic@orbimuinit\endcsname{perfnone}
\expandafter\def\csname blockcolour@dynamic@orbimuref\endcsname{perfnone}
\expandafter\def\csname blockcolour@dynamic@orbplacerec\endcsname{perfpoor}
\expandafter\def\csname blockcolour@dynamic@orbloopfusion\endcsname{perfnone}
\expandafter\def\csname blockcolour@dynamic@orbfullba\endcsname{perfnone}
\expandafter\def\csname blockcolour@dynamic@orbyolo\endcsname{perfnone}
\expandafter\def\csname blockcolour@dynamic@vinstrack\endcsname{perfgood}
\expandafter\def\csname blockcolour@dynamic@vinsqueue\endcsname{perfpoor}
\expandafter\def\csname blockcolour@dynamic@vinsimuprop\endcsname{perfgood}
\expandafter\def\csname blockcolour@dynamic@vinstri\endcsname{perfgood}
\expandafter\def\csname blockcolour@dynamic@vinsconstruct\endcsname{perfgood}
\expandafter\def\csname blockcolour@dynamic@vinssolve\endcsname{perfpoor}
\expandafter\def\csname blockcolour@dynamic@vinsmarg\endcsname{perfgood}
\expandafter\def\csname blockcolour@dynamic@vinsoutlier\endcsname{perfgood}
\expandafter\def\csname blockcolour@dynamic@vinsslide\endcsname{perfgood}
\expandafter\def\csname blockcolour@dynamic@vinsbackendtotal\endcsname{perfgood}
\expandafter\def\csname blockcolour@dynamic@vinssolverq\endcsname{perfgood}
\expandafter\def\csname blockcolour@dynamic@vinsloop\endcsname{perfnone}
\expandafter\def\csname blockcolour@dynamic@vinsrtab\endcsname{perfnone}
\expandafter\def\csname blockcolour@dynamic@vinsyolo\endcsname{perfnone}
%
  \resizebox{\textwidth}{!}{\pipelinefig{dynamic}}
  \caption[Per-block performance in the \textbf{dynamic} condition, over the twelve relogged dynamic runs of each estimator.]{Per-block performance in the \textbf{dynamic} condition, over the
  twelve relogged dynamic runs of each estimator. Loop closing carries no
  ablation here because the relogged campaign has no loop-closure-disabled
  variant. Evidence in \cref{tab:blockperf-dynamic-orb,tab:blockperf-dynamic-vins},
  one table per algorithm.}
  \label{fig:blockperf-dynamic}
\end{figure}
\clearpage
\begin{table}[p]
\centering
\caption[Per-block evidence for \textbf{ORB-SLAM3} in the dynamic condition, \cref{fig:blockperf-dynamic}.]{Per-block evidence for \textbf{ORB-SLAM3} in the dynamic condition, \cref{fig:blockperf-dynamic}. Twelve relogged dynamic runs, median across runs.}
\label{tab:blockperf-dynamic-orb}
\scriptsize
\setlength{\tabcolsep}{3pt}
%% Generated by script/block_performance.py -- do not edit.
\begin{tabularx}{\textwidth}{>{\raggedright\arraybackslash}p{0.155\textwidth}l>{\raggedright\arraybackslash}X r>{\raggedright\arraybackslash}p{0.20\textwidth}l}
\toprule
Block & Rule & Metric & Value & Criterion / note & Verdict \\
\midrule
\multicolumn{6}{l}{\textbf{Tracking}} \\
\addlinespace[2pt]
\quad Extract ORB (+preproc, stereo) & latency & median image\_preprocessing\_feature\_stereo\_ms & \num{45.1} & \textit{convention:} $<$40\% of stage (73.8\% of stage) & \perfmark{poor} \\
\addlinespace[3.5pt]
\quad IMU preintegration & latency & median imu\_preintegration\_ms & \num{0.116} & \textit{convention:} $<$40\% of stage (0.2\% of stage) & \perfmark{good} \\
\addlinespace[3.5pt]
\quad Initial pose estimation & latency & median pose\_prediction\_ms & \num{0.055} & \textit{convention:} $<$40\% of stage (0.1\% of stage) & \perfmark{good} \\
\addlinespace[3.5pt]
\quad Track local map & latency & median local\_map\_tracking\_ms & \num{12.2} & \textit{convention:} $<$40\% of stage (19.9\% of stage) & \perfmark{good} \\
\addlinespace[3.5pt]
\quad New keyframe decision & latency & median keyframe\_decision\_ms & \num{0.179} & \textit{convention:} $<$40\% of stage (0.3\% of stage) & \perfmark{good} \\
\addlinespace[3.5pt]
\quad TRACKING total & latency & median tracking\_total\_ms & \num{61.2} & \textit{derived:} $<$\SI{66.0}{\milli\second} frame budget & \perfmark{good} \\
\addlinespace[3.5pt]
\quad Dynamic feature removal (YOLO) & none & -- & -- & -- (filter disabled in every run reported here) & \perfmark{none} \\
\addlinespace[3.5pt]
\addlinespace[5pt]
\multicolumn{6}{l}{\textbf{Local mapping}} \\
\addlinespace[2pt]
\quad KeyFrame insertion & latency & median keyframe\_insertion\_ms & \num{10.9} & \textit{convention:} $<$40\% of stage (5.7\% of stage) & \perfmark{good} \\
\addlinespace[3.5pt]
\quad Recent MapPoints culling & latency & median map\_point\_culling\_ms & \num{0.714} & \textit{convention:} $<$40\% of stage (0.4\% of stage) & \perfmark{good} \\
\addlinespace[3.5pt]
\quad New points creation & latency & median map\_point\_creation\_fusion\_ms & \num{50.0} & \textit{convention:} $<$40\% of stage (26.3\% of stage) & \perfmark{good} \\
\addlinespace[3.5pt]
\quad Local BA & latency & median local\_ba\_ms & \num{122.9} & \textit{convention:} $<$40\% of stage (64.6\% of stage) & \perfmark{poor} \\
\addlinespace[3.5pt]
\quad Local keyframes culling & latency & median keyframe\_culling\_ms & \num{4.0} & \textit{convention:} $<$40\% of stage (2.1\% of stage) & \perfmark{good} \\
\addlinespace[3.5pt]
\quad IMU initialization & latency & median imu\_initialization\_ms & -- & \textit{convention:} $<$40\% of stage (not logged) & \perfmark{none} \\
\addlinespace[3.5pt]
\quad IMU scale refinement & latency & median inertial\_refinement\_ms & -- & \textit{convention:} $<$40\% of stage (not logged) & \perfmark{none} \\
\addlinespace[3.5pt]
\addlinespace[5pt]
\multicolumn{6}{l}{\textbf{Loop closing}} \\
\addlinespace[2pt]
\quad Place recognition & outcome & candidates accepted & 0 of 2 & \textit{convention:} $>$\SI{10}{\percent} of candidates accepted (5726 events) & \perfmark{poor} \\
\addlinespace[3.5pt]
\quad Loop fusion + essential graph & ablation & ATE(on)/ATE(off) & -- & \textit{convention:} $<$1.0 (enabling it should help) (one run per configuration) & \perfmark{none} \\
\addlinespace[3.5pt]
\quad Full BA & outcome & global BA executions & -- & -- (see \cref{tab:orb-loop}) & \perfmark{none} \\
\addlinespace[3.5pt]
\bottomrule
\end{tabularx}

\end{table}
\clearpage
\begin{table}[p]
\centering
\caption[Per-block evidence for \textbf{VINS-Fusion} in the dynamic condition, \cref{fig:blockperf-dynamic}.]{Per-block evidence for \textbf{VINS-Fusion} in the dynamic condition, \cref{fig:blockperf-dynamic}. Twelve relogged dynamic runs, median across runs.}
\label{tab:blockperf-dynamic-vins}
\scriptsize
\setlength{\tabcolsep}{3pt}
%% Generated by script/block_performance.py -- do not edit.
\begin{tabularx}{\textwidth}{>{\raggedright\arraybackslash}p{0.155\textwidth}l>{\raggedright\arraybackslash}X r>{\raggedright\arraybackslash}p{0.20\textwidth}l}
\toprule
Block & Rule & Metric & Value & Criterion / note & Verdict \\
\midrule
\multicolumn{6}{l}{\textbf{Front end}} \\
\addlinespace[2pt]
\quad Feature tracking & latency & median feature\_tracking\_ms & \num{42.5} & $<$\SI{66.0}{\milli\second} frame budget & \perfmark{good} \\
\addlinespace[3.5pt]
\quad Feature queue & latency & median feature\_queue\_wait\_ms & \num{502.1} & $<$\SI{66.0}{\milli\second} frame budget (backlog 7.5 frames) & \perfmark{poor} \\
\addlinespace[3.5pt]
\quad Persistent-feature removal (YOLO) & none & -- & -- & -- (filter disabled in every run reported here) & \perfmark{none} \\
\addlinespace[3.5pt]
\addlinespace[5pt]
\multicolumn{6}{l}{\textbf{Back end}} \\
\addlinespace[2pt]
\quad IMU propagation & latency & median imu\_propagation\_ms & \num{0.204} & \textit{convention:} $<$40\% of stage (0.3\% of stage) & \perfmark{good} \\
\addlinespace[3.5pt]
\quad Triangulation & latency & median triangulation\_ms & \num{0.062} & \textit{convention:} $<$40\% of stage (0.1\% of stage) & \perfmark{good} \\
\addlinespace[3.5pt]
\quad Problem construction & latency & median problem\_construction\_ms & \num{1.1} & \textit{convention:} $<$40\% of stage (1.7\% of stage) & \perfmark{good} \\
\addlinespace[3.5pt]
\quad Ceres solve & latency & median solver\_ms & \num{57.2} & \textit{convention:} $<$40\% of stage (89.8\% of stage) & \perfmark{poor} \\
\addlinespace[3.5pt]
\quad Marginalization & latency & median marginalization\_ms & \num{2.4} & \textit{convention:} $<$40\% of stage (3.7\% of stage) & \perfmark{good} \\
\addlinespace[3.5pt]
\quad Outlier rejection & latency & median outlier\_rejection\_ms & \num{0.068} & \textit{convention:} $<$40\% of stage (0.1\% of stage) & \perfmark{good} \\
\addlinespace[3.5pt]
\quad Slide window & latency & median slide\_window\_ms & \num{0.167} & \textit{convention:} $<$40\% of stage (0.3\% of stage) & \perfmark{good} \\
\addlinespace[3.5pt]
\quad BACK END total & latency & median backend\_total\_ms & \num{63.7} & $<$\SI{66.0}{\milli\second} frame budget & \perfmark{good} \\
\addlinespace[3.5pt]
\quad Solver convergence & outcome & solves returning a usable solution & 99.55\% & \textit{convention:} $\geq$\SI{99}{\percent} usable (\SI{45.1}{\percent} hit the iteration or \SI{80}{\milli\second} cap without converging; reported, not graded) & \perfmark{good} \\
\addlinespace[3.5pt]
\addlinespace[5pt]
\multicolumn{6}{l}{\textbf{Loop closing}} \\
\addlinespace[2pt]
\quad Loop fusion & none & -- & -- & -- (not implemented in this fork) & \perfmark{none} \\
\addlinespace[3.5pt]
\quad RTAB-Map correction & none & -- & -- & -- (separate process; correction never injected) & \perfmark{none} \\
\addlinespace[3.5pt]
\bottomrule
\end{tabularx}

\end{table}
\clearpage


Four readings follow that the trajectory tables alone do not give.

\textbf{Both visual front ends are the dominant cost, and neither keeps up
cleanly.} ORB-SLAM3 spends \SI{33.1}{\milli\second} of a \SI{53.8}{\milli\second}
tracking budget on preprocessing, extraction and stereo matching in static scenes,
rising to \SI{45.1}{\milli\second} of \SI{61.2}{\milli\second} in dynamic ones.
VINS-Fusion's back end waits \SI{403.7}{\milli\second} per frame on its feature
queue in static scenes and \SI{502.1}{\milli\second} in dynamic ones, against a
\SI{66.0}{\milli\second} frame budget, with a backlog of four to seven frames.
\flag{These are the same deficit wearing different clothes}: neither estimator
processes frames as fast as they arrive, and the two differ in what they do about
it. ORB-SLAM3 subscribes best-effort at depth five and discards the overflow,
which is the frame shortfall of \cref{sec:results-validity}; VINS-Fusion
subscribes reliably at depth 100 and absorbs it as latency. The frame counters see
one and not the other.

\textbf{Local BA dominates ORB-SLAM3's local mapping} at \SI{244.4}{\milli\second}
of \SI{312.9}{\milli\second} in static scenes. It runs per keyframe rather than
per frame, so it does not compete directly with the frame budget, but it is where
that pipeline's compute goes.

\textbf{Loop closing is not earning its place in these scenes.} Across the five
\repo{allsensor} datasets, enabling it leaves the median aligned ATE
\num{1.37} times \emph{worse} than disabling it, and only 4 of 281 place-recognition
candidates were ever accepted (\cref{tab:orb-loop}). This is the project's one
ORB-SLAM3 ablation, and it is still one run per configuration, so it sits inside
an unmeasured run-to-run spread---the verdict is reported, not relied upon.

\textbf{The EKF's accuracy rests on one intervention.} Removing the wheel update
costs a factor of \num{54} in aligned ATE; removing the IMU prediction costs
\num{6.6}. The \num{54} covers the wheel speed correction and the gyroscope-bias
observation \emph{together}: both are measurement rows of the same update, one
flag removes both, and no flag removes only one, so
\cref{tab:blockperf-static-ekf} lists the figure against each row while marking it
as joint rather than presenting two independent results. Swapping the yaw-rate model from the bicycle formulation to the
rear-wheel differential, by contrast, changes nothing measurable---median
\num{1.0}, and on \repo{allsensor_004} the noisier estimator is the better one.
That is worth stating plainly, because the differential model is four times worse
as a \emph{rate} estimator (\cref{tab:proprio-noise}) and dominates standalone
wheel odometry's error. Inside the filter it barely matters, because the wheel
rate is used only to observe the gyroscope bias and is weighted by its own
measured $\sigma$. The design choice of \cref{sec:proprio-method} is vindicated by
its own ablation.

The campaign succeeds as an instrumentation milestone: repeated run-level data,
front-end and back-end timing, input retention, process resources, initialization,
tracking loss, keyframes, loop closures, map merges, and raw trajectories are now
available in auditable form. Ground-truth accuracy can now be measured, but the
campaign does not support the intended semantic-filter hypothesis or a clean
dynamic-scene estimator ranking. The limitations are:

\Cref{tab:limitations} lists them with the consequence each one carries.

\begin{table}[htbp]
\centering
\caption[Limitations of the results reported in this chapter, and what each one prevents.]{Limitations of the results reported in this chapter, and what each one
prevents. ``Scope'' identifies the experiment affected: R is the relogged
campaign, E the estimator comparison of \cref{sec:results-proprio}, and A both.}
\label{tab:limitations}
\small
\begin{tabularx}{\textwidth}{l>{\raggedright\arraybackslash}p{0.36\textwidth}>{\raggedright\arraybackslash}X}
\toprule
Scope & Limitation & Consequence \\
\midrule
A & RPE was missing & now implemented and reported (\cref{sec:rpe}); it
  withdrew this chapter's local-coherence reading \\
R & every dynamic trajectory fails the jump-plausibility criterion & no dynamic
  estimator ranking \\
R & the VINS semantic treatment was inactive in all repeats & no VINS filtering
  conclusion \\
R & 11 of 12 ORB treatment pairs fail the paired-input rule & no cross-treatment
  accuracy conclusion \\
R & YOLO truth is projected simulator geometry, not annotated visible extent, and
  the retained confidence threshold is fixed & no AP claim \\
R & class metrics are conditional on successful IoU localization & must not be
  read as end-to-end accuracy \\
R & two static ORB run summaries have incomplete schemas & retained as
  operational; counters unavailable \\
A & CPU and memory are per-estimator-process figures & not whole-system cost \\
A & one laptop and one replay rate throughout & no hardware or input-rate stress
  surface \\
E & \flag{replay rate $0.5\times$ for both visual estimators} & accuracy figures
  are not evidence of real-time capability \\
E & \flag{baselines are offline, visual estimators ran live} & timing-induced
  error is present on one side of the comparison only \\
E & \flag{only \repo{allsensor} bags carry wheel encoders} & the baseline exists
  for static featureless-floor scenes only; no dynamic condition \\
E & \flag{$n=5$ across four worlds, no condition repeated more than twice} & no
  mean, deviation or interval; between-bag differences are not effects \\
E & \flag{ORB-SLAM3 drops frames at a depth-5 best-effort subscriber; VINS buffers
  100 deep} & input delivery is not matched between the two; no timing comparison
  is admissible \\
E & \flag{ORB-SLAM3 run-to-run spread unmeasured, one run per configuration} &
  loop-closure differences cannot be attributed; repeats required \\
E & \flag{loop-closure state absent from \repo{run_metadata.csv}} & configurations
  distinguishable only by inspecting \repo{loop_closing.csv} \\
E & \flag{origin of the wheel yaw-rate offset untested} & the drift is predicted
  but not yet attributable or correctable (\cref{sec:h1-falsified}) \\
E & wheel-rate offset is absorbed by the EKF gyroscope-bias state & now
  quantified: EKF heading error rises tenfold where the wheel rate is poor
  (\cref{tab:ekf-coupling}) \\
E & \flag{accelerometer characterisation unstable across bags} & bias 0.019 to
  0.068, scale 1.020 to 1.067 for one simulated sensor \\
E & neither visual estimator writes a pose covariance & the consistency test
  cannot be applied to them \\
\bottomrule
\end{tabularx}
\end{table}

No root cause is assigned in this report. In particular, differences are not
attributed to calibration, timestamp handling, masks, replay scheduling, or a
specific estimator component without a controlled diagnostic experiment. The
immediate advisor decision is whether the next work period should first restore
active, input-matched filtering and plausible full trajectories, or
instead prioritize the planned integration of relocalization, global
optimization, and map merging. The latter remains future work and would provide
a unified localization/map input for navigation; path planning itself remains
outside scope.
